\section{Análisis de casos}

\subsection{Soang: de la búsqueda a DPO}

Soang es un juego de box-pushing en el que solo se puede empujar hacia adelante, y pone a prueba en gran medida la capacidad del sistema para planificar los estados futuros. El enfoque de Compound AI se divide en dos etapas: primero se usa el trace de búsqueda para aprender el orden estricto de las acciones, y después se usa DPO para que el agent judge dé retroalimentación y reduzca el action spirit. El flujo de entrenamiento es el siguiente:
\begin{enumerate}
    \item Se usa el solver para generar el trace completo, registrando cost, plan e intermediate state;
    \item Se etiqueta el trace (si cumple las constraints, si se completa en 4 rondas);
    \item Search Former aprende el trace; el agent judge muestrea varios plan y evalúa la fidelity;
    \item El fine-tuning con DPO hace que el agent judge se incline más hacia el plan óptimo.
\end{enumerate}
El resultado final es que Search Former + DU Forer, con 15M de parámetros y 100K trace, alcanza una tasa de éxito del 80\%, mientras que el baseline solution-only necesita 175M de parámetros y 1M de trace para llegar solo al 20\%.
\begin{table}[H]
\centering
\begin{tabular}{lrrl}
\toprule
Modelo & Parámetros & Muestras & Tasa de éxito \\
\midrule
Solution-only & 175M & 1M & 20\% \\
Search Former + DU Forer & 15M & 100K & 80\% \\
\bottomrule
\end{tabular}
\caption{Comparación de modelos en el benchmark de Soang}
\end{table}

\subsection{Nano optics: del latent cost a la fabricación}

En el diseño óptico de una malla de 80\u00d780, el trace de búsqueda se convierte en un trace de simulación: el agent observa la trayectoria de la luz, la respuesta en frecuencia y el interference pattern, y después comprime esta información en el latent cost \(C\). Todo el pipeline incluye: 1) generar la simulación del beam splitter y el wavelength multiplexer; 2) elegir el binary pattern que satisface la restricción de brush; 3) usar el solver para producir \(X^\star\); 4) propagar el loss hacia atrás hasta \(C\).
\begin{importantbox}{El ciclo cerrado del trace a la fabricación}
1) El trace conserva la información de la interferencia de la luz y de la brocha de fabricación; 2) el latent \(C\) extrae la descripción del entorno y a la vez controla al solver; 3) la retroalimentación del solver se usa para calibrar \(C\) y evitar generar puntos no fabricables; 4) el multi-task loss permite que el diseño se transfiera entre distintos benchmark.
\end{importantbox}
\begin{knowledgebox}{Enseñanzas de los casos}
Estos dos casos muestran en conjunto que Compound AI puede encadenar el trace de planificación, las preguntas y respuestas interactivas y el Latent solver, de modo que cada salida es a la vez la entrada de la siguiente etapa y un resultado intermedio auditable. Este modular pipeline facilita localizar las fallas y hacer iteraciones en pasos pequeños.
\end{knowledgebox}

\subsection{Resumen de la sección}

El análisis de casos destaca el flujo de datos en la práctica de la ingeniería: el search trace, el agent judge y el latent cost se conectan entre sí, y se validan repetidamente en distintas tareas para garantizar que la teoría abstracta pueda llevarse a la práctica.

